\documentclass[11pt,a4paper]{article}
\usepackage[margin=2.2cm]{geometry}
\usepackage[T1]{fontenc}
\usepackage{amsmath,amssymb,xcolor,enumitem,booktabs,graphicx,tabularx}
\usepackage[normalem]{ulem}
\setlength{\parskip}{0.6em}\setlength{\parindent}{0pt}
\definecolor{newc}{HTML}{1f6f43}
\newcommand{\new}[1]{\textcolor{newc}{#1}}
\newcommand{\old}[1]{\textcolor{gray}{\sout{#1}}}
\newcommand{\note}[1]{{\small\itshape [#1]}}
\newcommand{\lab}[1]{\textsf{\footnotesize\detokenize{#1}}}
\renewcommand{\ref}[1]{\lab{#1}}
\renewcommand{\eqref}[1]{(\lab{#1})}
\newcommand{\citep}[1]{[\lab{#1}]}
\newcommand{\citet}[1]{\lab{#1}}
\title{Proposed change, Introduction, Project Description,\\ Background and Related Work}
\date{24 November 2026, version 1}
\begin{document}\sloppy
\maketitle

\section*{How to read this}
This answers sections 3 and 4 of \texttt{Checklist.md}, as corrected on
24 November. The aim is to stop telling the same thing several times.
Agreed on 24 November: each part gets one job. The Introduction gives the
problem, the coupling and what is compared. Aim and scope gives the factors
and the scope. The research questions give only the questions. Expected
outcome gives only the predictions. The tragedy of the commons stays,
because it returns in the Discussion.

Black text stays, \new{green} is new and \old{struck} text goes.
\note{Notes in brackets} say why. Nothing below is in the thesis yet.
Parts~1 and~2 show the whole text. Parts~3 and~4 show only the paragraphs
that change.

\section*{What was repeated}
\begin{itemize}[nosep]
  \item ``The whole fleet, one zone or one robot'' was told four times
    (Introduction, Aim and scope, the main question, RQ1). Now once, in the
    Introduction.
  \item ``The main load and a higher load where the aisles fill'' was told
    four times. Now the Introduction gives the reason, RQ2 asks it and
    Expected outcome predicts it.
  \item The popular SKU near a workstation was the example of two paragraphs
    in a row. Now one.
  \item The percentage gain paragraph after the questions is nearly word for
    word the one beside \lab{eq:relgain} in the Problem Formulation.
  \item Orders, pick tasks and put-away tasks were explained in Aim and scope
    and again, properly, in the Problem Formulation.
  \item ``Three problems usually solved apart'' opens the Contributions, the
    Problem Formulation and the research gap.
\end{itemize}

% ================================================================
\section*{Part 1. Introduction}

An automated warehouse makes two kinds of decision that depend on each
other. The storage layer decides where stock-keeping units (SKUs) are kept.
The transport layer gives pickup and delivery tasks to robots and routes
them through shared aisles. The first is studied as the Storage Location
Assignment Problem (SLAP), including dynamic versions where the placement
changes over time. The second is studied as Multi-Agent Path Finding (MAPF)
and Multi-Agent Pickup and Delivery (MAPD) \citep{ma2017lifelong_mapf},
where many robots keep carrying out tasks without colliding.

Storage and routing affect each other. Where the stock sits decides where
the robots must go. \old{If a popular SKU is put on a shelf near a
workstation, many trips go to that shelf and its aisle gets busy.} In the
other direction, the traffic decides how long a shelf really takes to
reach. \old{A shelf next to a workstation is still slow to reach if its
aisle is full of robots.}
\old{So neither layer can be planned well on its own.} Suppose storage puts
every popular SKU on the shelves closest to the workstations. On paper this
gives the shortest trips. In practice every robot now goes to the same few
aisles, they block each other, and orders take longer than before.
\new{So neither layer can be planned well on its own.}
\note{3.1. Two paragraphs with the same example become one. Both directions
stay, and the example is told once. The conclusion moves to the end, where
the example has earned it.}

This thesis measures how much it helps to move stock as demand changes, and
whether that gain depends on how the robots are controlled. A storage rule
decides from time to time where arriving stock goes, based on the orders it
has seen. Between these updates, a controller gives tasks to the robots and
routes them. Demand storage is compared with fixed storage under three
controllers. The centralised controller decides for the whole fleet, the
section-based one for one zone and the decentralised one for a single
robot. All three are graph neural networks trained the same way with
reinforcement learning, so they differ only in how much one decision
covers. Token Passing, a classic planner that does not learn, is shown next
to them as a reference.
\note{Unchanged. This is now the one place that says what each controller
decides for.}

The coupling is tested in two ways. The first runs one way. Demand storage
changes where the robots go, and the thesis asks whether the gain it gives
depends on the controller. If it does, what storage is worth depends on how
the robots are routed.

The second runs both ways. The storage rule can also read where robots get
stuck, and steer both new stock and picks away from those shelves. Then the
robots' movement feeds back into where the stock goes and where it is taken
from, and the loop between the two layers is closed. Avoiding crowded
aisles can only help when the aisles are crowded. So the loop is tested at
the main load and again at a higher load\old{, where the fleet still keeps
up but the aisles fill}.
\note{The detail of the higher load is in RQ2.}

\new{The decentralised robots may also send short messages to the robots
near them. A robot that decides alone sees only what is around it, and a
message can tell it about crowding further away.}
\note{New. RQ4 asks about communication, but the Introduction never
mentioned it. See question~1.}

The main measure is the mean time an order takes, from when it is placed
until it leaves the warehouse.

% ================================================================
\section*{Part 2. Project Description}

\subsection*{Aim and scope}

The project aims to find out how feedback between storage and the control
of the robots affects how well a warehouse works. \old{The warehouse is a
finite graph. Customers keep sending orders. An order is a short list of
items. Each item becomes a pick task, where a robot fetches it from its
shelf and brings it to a workstation. So where the stock sits decides where
the robots go. New stock also arrives at the warehouse. Each delivery
becomes a put-away task, where a robot carries it to a shelf chosen by the
storage rule.}
\note{3.2. The Problem Formulation defines orders and tasks properly, and
the Introduction already says that the stock decides where the robots go.}

Four things are varied in the experiments:
\begin{enumerate}[nosep]
    \item the storage rule: fixed, demand or congestion-aware
    \item the controller: centralised, section-based or decentralised\old{,
      each a learned policy that decides for the whole fleet, one zone or one
      robot}
    \item communication: off or local, for the decentralised controller
      only, because the other two already see more than one robot
    \item the load: the number of robots and how fast orders arrive, as a
      share of what the fleet can serve
\end{enumerate}
The main result compares demand storage with fixed storage under each
controller, at a main load where every controller should keep up. The
other factors test how far that result reaches. \old{Congestion-aware
storage tests whether closing the loop helps. It is run at the main load
and at a higher load, where the aisles fill. A sweep of rising loads then
shows how busy the warehouse can get before the gain of demand storage
disappears.}
\note{RQ2 and RQ3 ask exactly this.}

\old{The mean order service time gives the main result. It is split into the
wait for a robot, travel and blocked time, which shows which part the
storage rule changes. Blocked time and how crowded the aisles get show
whether closing the loop helps.} \new{A fleet keeps up when it} \old{Whether
the fleet keeps up, meaning it} finishes tasks as fast as they arrive\old{,
shows where the gain stops}. Order times are compared only when both runs
keep up, because a fleet that falls behind has a queue that grows for as
long as the run lasts. \old{Where they cannot be compared, the share of the
arriving work each fleet serves is compared instead. That share also shows
whether messages help the decentralised controller.}
\note{3.2. The split of service time and the served share are set out in
the Experiments chapter. Only ``keeps up'' stays, because RQ1 and RQ2 use
it.}

\new{The study is a simulation of identical robots. It leaves out people,
batteries and how a robot drives. Stock already on a shelf never moves, so
the cost of moving it does not arise. The Problem Formulation lists the
full scope.}
\note{New. The Checklist asks Aim and scope to say what is outside the
study. Taken from \lab{sec:pf:scope}. See question~2.}

\subsection*{Research question}

The main research question is: \newline
\textbf{RQ:} How much does demand storage shorten the service time of
orders, how does that gain differ between \new{the three controllers}
\old{controllers that decide for the whole fleet, one zone or one robot},
and does closing the loop add to it?

It is split into four questions:
\begin{description}[nosep]
  \item[RQ1:] Does demand storage shorten order service time compared with
  fixed storage, at a load where the fleet keeps up, and how does that gain
  differ between \new{the three controllers}\old{controllers that decide for
  the whole fleet, one zone or one robot}?
  \item[RQ2:] Does closing the loop, so that storage also avoids shelves
  where robots get blocked, make orders faster than demand storage does,
  at the main load and at a higher load where the aisles are crowded?
  Where a fleet falls behind, does it let the fleet serve more?
  \item[RQ3:] As orders arrive faster, does the gain of RQ1 grow, shrink or
  disappear, and at what rate does each trained controller stop keeping up?
  \item[RQ4:] Does local communication make the decentralised controller
  route better than what a robot sees on its own?
\end{description}

\old{The gain of RQ1 is given as a percentage of the time orders take under
fixed storage. A controller that is slow under fixed storage has more to
save, so the gain in steps alone would favour the controller that routes
worst.}
\note{3.3. The same reason, with its citation, stands beside
\lab{eq:relgain} in the Problem Formulation, where the gain is defined.}

\subsection*{Contributions}

\old{The thesis joins three problems that are usually solved apart, storage
assignment, task assignment and path finding.} \new{The thesis} \old{It}
makes \new{three}\old{four} contributions.
\begin{enumerate}[nosep]
\item One model in which the three problems share the same warehouse graph
  and the same stock, so each changes what the others face
  (Chapter~\ref{sec:problem}).
\item Congestion-aware storage, the thesis's own rule. It adds the time
  picks stand still at a shelf to the classic slotting score, so where new
  stock goes reacts to the traffic of the agents
  (Section~\ref{sec:method:storage}).
\item A simulator built on that model. Between two runs only the storage
  rule or only the controller changes. A difference between fixed and demand
  storage can then be traced to storage itself and not to a different setup.
\item \old{Results that answer each research question, for the warehouses,
  demand and controllers used in the experiments
  (Chapter~\ref{sec:experiments}).}
\end{enumerate}
\note{3.4. The lead-in is told again in the Problem Formulation opening and
the research gap, and item~1 says it more precisely. Item~4 is what every
thesis delivers, not a contribution. See question~3.}

\subsection*{Expected outcome}
\note{Unchanged. It holds only the predictions, as agreed.}

% ================================================================
\section*{Part 3. Background, the paragraphs that change}

\subsection*{Multi-Agent Path Finding}

\ldots This is fast, but it promises neither the cheapest paths nor a path
for every agent.

Conflict-Based Search (CBS) sits between the two. It \new{plans each agent
alone and plans again only for the agents whose paths
clash~\citep{sharon2015cbs}.} \old{first plans each agent as if the others
were absent. Where two plans clash, it tries two fixes, each forbidding one
of the two agents from that place at that step, and plans again only for
that agent. It searches these fixes until it finds the cheapest set of
paths with no clash~\citep{sharon2015cbs}.}
\note{4.2. No later chapter uses CBS. One sentence keeps the picture of the
solver families.}

\subsection*{Reinforcement learning}

\ldots the actor learns from how much better or worse an action turned out
than the critic expected. \old{Not every method has a critic.} Deep
Q-learning\old{, for example,} \new{has no critic. It} learns the value of
each action directly and acts on the best one.
\note{4.3. Kept, because Related Work uses Deep Q-learning twice (van
Knippenberg, Ma 2021). Only the framing goes. It has no citation today, see
question~4.}

\old{In methods with a critic, what the critic may see is a design choice.
In centralised training with decentralised execution, the critic sees the
whole state during training, while each agent still acts only on what it
sees~\citep{lowe2017maddpg}.} \new{Agents may also} \old{A different line of
work lets agents} choose their actions one after another, each seeing the
choices before it. Any joint choice can be built this way, while the
network only runs once per agent~\citep{wen2022mat,fu2022revisiting}.
\note{4.3. The Method explains centralised training with its citation where
it is rejected ([A128]). Choosing one after another stays, because SePar
and the Method use it.}

\subsection*{Graph neural networks}

\ldots The second kind is how agents send each other messages.

\old{GNNs have been combined with reinforcement learning for routing on
graphs. \citet{tien2022vrp_drl_gnn} use a GNN with deep reinforcement
learning to solve vehicle routing problems. \citet{lischka2024great} build
a GNN around the edges for problems such as the travelling salesman
problem, and argue that edge features carry information that a network
over vertices alone loses.}
\note{4.4. Nothing later cites these, and the Tien link is broken. See
question~5 on Lischka.}

\subsection*{Congestion}
\note{Unchanged. The tragedy of the commons stays for the Discussion.}

% ================================================================
\section*{Part 4. Related Work, the paragraph that changes}

\ldots \citet{li2021magat} extend this network with attention, so each robot
weighs the messages it receives. \old{Its success rate is up to 47\% above
the benchmark they compare with, also on}\new{It also does well on} problems
100 times larger than those it was trained on. \ldots
\note{4.6. A benchmark unlike this thesis's, and nothing uses the number.
The 6.3\% of Waubert de Puiseau et al.\ (the size of a storage gain), the
1024 agents of PRIMAL (scale) and the switch-off test of Pham et al.\ (it
bears on RQ4 and the crowd reward) stay.}

% ================================================================
\section*{Questions}
\begin{enumerate}[nosep]
  \item The new Introduction paragraph on messages. Keep it, or is RQ4 meant
    to stay minor and the Introduction silent about it?
  \item The scope paragraph in Aim and scope. Keep it as written, cut it to
    one sentence that points to the Problem Formulation, or leave it out?
  \item Contribution~4. Drop it now, or keep it until the results are in and
    then name what the results show?
  \item Deep Q-learning has no citation. Add Mnih et al.\ (2015) to
    \texttt{references.bib}, or leave it uncited as it is today?
  \item Lischka et al. Drop it, or cite it in the Method beside the move
    scores, since it argues that features on the edges help?
\end{enumerate}

\end{document}
